\documentclass[11pt,a4paper]{article}
\usepackage[utf8]{inputenc}
\usepackage[T1]{fontenc} % ESSENCIAL: Resolve a quebra de caracteres no ATS
\usepackage{lmodern} % ESSENCIAL: Garante fonte mapeável e selecionável no PDF
\usepackage[margin=0.8in]{geometry}
\usepackage{titlesec}
\usepackage{enumitem}
\usepackage[hidelinks]{hyperref} % hidelinks remove a borda ao redor dos links

\titleformat{\section}{\Large\bfseries}{}{0em}{}[\titlerule]
\titlespacing{\section}{0pt}{1.5ex}{1ex}
\setlength{\parindent}{0pt}
\setlist[itemize]{leftmargin=*, nosep, itemsep=0.2em}

\begin{document}
\pagestyle{empty}

{\centering
\textbf{\huge Marcio Gabriel Schinor Mazega} \\ \vspace{0.2cm}
\textit{\large Cientista de Dados | Machine Learning} \\ \vspace{0.1cm}
Tupã, SP -- Brasil | (14) 99812-6535 | \href{mailto:marciomazegaa@gmail.com}{marciomazegaa@gmail.com} \\
\href{https://linkedin.com/in/marcio-mazega}{linkedin.com/in/marcio-mazega} | \href{https://github.com/marciomazega}{github.com/marciomazega} \\
}

\vspace{0.4cm}
\section*{RESUMO PROFISSIONAL}
Cientista de Dados com forte viés em modelagem preditiva, processamento de linguagem natural (NLP) e inteligência geoespacial (AgTech). Desenvolvo análises complexas utilizando Python e Pandas, integrando algoritmos de Machine Learning (Scikit-learn) para descoberta de padrões e tomada de decisão. Possuo rigor científico validado por pesquisas (FAPESP/CNPq) e experiência corporativa arquitetando pipelines analíticos e validando modelos de Inteligência Artificial, aliando a fundamentação acadêmica à geração direta de valor para o negócio.

\vspace{0.2cm}
\section*{HABILIDADES TÉCNICAS}
\begin{itemize}
    \item \textbf{IA e Machine Learning:} Modelagem Preditiva, Scikit-learn, Algoritmos de Regressão/Classificação, NLP.
    \item \textbf{Análise de Dados:} Python, Pandas, NumPy, Análise Geoespacial (Google Earth Engine/GIS), Estatística.
    \item \textbf{Banco de Dados e Infraestrutura:} PostgreSQL, MySQL, Docker, Linux (Ubuntu/WSL2).
    \item \textbf{Ferramentas e Deploy:} Git / GitHub, FastAPI (Deploy de Modelos), Web Scraping.
\end{itemize}

\vspace{0.2cm}
\section*{EXPERIÊNCIA PROFISSIONAL}
\textbf{Engenheiro de Dados e Analytics} \hfill \textit{Fev 2026 -- Jun 2026} \\
\textit{Gattaz Health} \\
\begin{itemize}
    \item Construí e orquestrei pipelines analíticos para o processamento unificado de [X] milhões de registros clínicos.
    \item Tratei, higienizei e cruzei dados sensíveis em PostgreSQL, otimizando o consumo em modelos de BI em [X]\%.
    \item Apliquei técnicas de exploração e normalização de dados para garantir a integridade e a escalabilidade das análises corporativas.
\end{itemize}

\textbf{Validador de Inovação em IA (Projeto)} \hfill \textit{Mai 2026 -- Atual} \\
\textit{Move Strategy Consulting (Sentinela IA \& Paciente Zero Tracker)} \\
\begin{itemize}
    \item Validei a estrutura e a viabilidade técnica de projetos de Inteligência Artificial voltados ao monitoramento e saúde.
    \item Modelei a arquitetura lógica para sistemas de rastreamento preditivo (Paciente Zero), alinhando requisitos de dados às metas de inovação da companhia.
\end{itemize}

\vspace{0.2cm}
\textbf{Pesquisador Científico (Bolsista FAPESP)} \hfill \textit{Ago 2025 -- Fev 2026} \\
\textit{Centro Avançado de Pesquisa Smart B100} \\
\begin{itemize}
    \item Conduzi pesquisas voltadas a metodologias de Living Labs, analisando [X] bases de dados de sistemas digitais rurais.
    \item Mapeei e modelei soluções AgTech orientadas a dados, criando matrizes de competitividade para o aprimoramento do manejo inteligente.
\end{itemize}

\vspace{0.2cm}
\textbf{Monitor de Iniciação Científica (ICJr)} \hfill \textit{Jan 2025 -- Mai 2026} \\
\textit{OBA - Olimpíada Brasileira de Astronomia e Astronáutica} \\
\begin{itemize}
    \item Liderei a análise estatística de dados experimentais, garantindo a precisão metodológica nos projetos de [X] alunos.
\end{itemize}

\vspace{0.2cm}
\section*{PROJETOS DE DESTAQUE}
\textbf{Buscador Inteligente de Artigos Acadêmicos} | \href{https://github.com/marciomazega}{[Ver no GitHub]}
\begin{itemize}
    \item \textbf{Stack:} Python, NLP, FastAPI, Web Scraping.
    \item \textbf{Descrição:} Desenvolvi uma solução automatizada para realizar varredura, extração e indexação inteligente de mais de [X] publicações, utilizando processamento de linguagem natural (NLP) para otimizar fluxos de revisão bibliográfica.
\end{itemize}

\textbf{Central Agro -- Inteligência Agronômica Geoespacial}
\begin{itemize}
    \item \textbf{Stack:} Python, Google Earth Engine (GEE), Scikit-learn.
    \item \textbf{Descrição:} Criei um pipeline preditivo capaz de cruzar Índices de Vegetação (NDVI) extraídos de satélite com variáveis meteorológicas, gerando insights antecipados para mitigação de riscos agrícolas.
\end{itemize}

\textbf{PRF Collision Predictor -- Modelo de Risco Rodoviário} | \href{https://github.com/marciomazega}{[Ver no GitHub]}
\begin{itemize}
    \item \textbf{Stack:} Python, Pandas, Seaborn, Scikit-learn.
    \item \textbf{Descrição:} Treinei e avaliei modelos de Machine Learning de classificação e regressão para prever a severidade de acidentes rodoviários com base em padrões espaciais e temporais históricos.
\end{itemize}

\vspace{0.2cm}
\section*{FORMAÇÃO ACADÊMICA E CERTIFICAÇÕES}
\textbf{Graduação em Tecnologia em Sistemas Inteligentes (TSI)} \hfill \textit{Jan 2025 -- Dez 2027} \\
\textit{FATEC Pompéia "Shunji Nishimura"} \\
\vspace{0.1cm}
\textbf{Certificações em Destaque:}
\begin{itemize}
    \item Introdução ao Machine Learning -- MBA USP/ESALQ (2026)
    \item Banco de Dados Relacional -- EncData
\end{itemize}

\vspace{0.2cm}
\section*{PREMIAÇÕES E PUBLICAÇÕES}
\begin{itemize}
    \item \textbf{Premiação:} 3º Lugar -- DSIN Coder Challenge (Desafio estadual de lógica e programação, 2025).
    \item \textbf{Publicações:} Coautor de 2 capítulos de livros acadêmicos sobre Inovação no Agronegócio e Sociologia.
\end{itemize}

\end{document}